% !TeX program = lualatex
% Compilacion recomendada:
% latexmk -lualatex Informe_datos_casos_estudio_papers.tex
\RequirePackage{fix-cm}
\documentclass[11pt,a4paper,oneside]{report}

\usepackage{fontspec}
\setmainfont{Arial}
\setsansfont{Arial}
\usepackage[spanish,es-nodecimaldot,es-noshorthands]{babel}
\usepackage{geometry}
\geometry{
  a4paper,
  left=3cm,
  top=3cm,
  right=2.5cm,
  bottom=2.5cm,
  headheight=14pt,
  headsep=0.6cm
}
\usepackage{setspace}
\onehalfspacing
\usepackage{microtype}
\usepackage{amsmath}
\usepackage{amssymb}
\usepackage{booktabs}
\usepackage{longtable}
\usepackage{array}
\usepackage{calc}
\usepackage{ragged2e}
\usepackage{enumitem}
\usepackage{etoolbox}
\usepackage{titlesec}
\usepackage{fancyhdr}
\usepackage{xcolor}
\usepackage{colortbl}
\usepackage[hidelinks]{hyperref}
\usepackage{bookmark}
\newcounter{none}

\definecolor{unired}{HTML}{8B1A1A}
\definecolor{darkblue}{HTML}{17365D}
\definecolor{lightgray}{HTML}{F2F2F2}

\hypersetup{
  pdftitle={Informe de datos de casos de estudio y referencias de verificación térmica},
  pdfauthor={Luis Enrique Koc Góngora y Herbert Antonio Meléndez García},
  pdfsubject={Ejemplos de datos recolectados del objeto de estudio}
}

\newcommand{\tesisbodyfont}{\fontsize{11}{13.6}\selectfont}
\setlength{\parindent}{0pt}
\setlength{\parskip}{12pt}
\setlength{\emergencystretch}{3em}
\renewcommand{\arraystretch}{1.18}
\setlength{\LTpre}{0.45\baselineskip}
\setlength{\LTpost}{0.65\baselineskip}
\AtBeginEnvironment{longtable}{%
  \footnotesize
  \setlength{\tabcolsep}{3pt}%
  \setlength{\parindent}{0pt}%
}
\setlist{
  leftmargin=1.1cm,
  before=\begin{spacing}{1}\tesisbodyfont,
  after=\end{spacing},
  topsep=.5\baselineskip,
  itemsep=.35\baselineskip,
  parsep=0pt,
  partopsep=0pt
}

\providecommand{\tightlist}{%
  \setlength{\itemsep}{0pt}\setlength{\parskip}{0pt}}

\titleformat{\section}
  {\normalfont\bfseries\fontsize{11}{13.2}\selectfont}
  {\thesection}{0.75em}{}
\titlespacing*{\section}{0pt}{1.0\baselineskip}{0.25\baselineskip}
\titleformat{\subsection}
  {\normalfont\bfseries\fontsize{11}{13.2}\selectfont}
  {\thesubsection}{0.75em}{}
\titlespacing*{\subsection}{0pt}{0.75\baselineskip}{0.25\baselineskip}
\titleformat{\subsubsection}
  {\normalfont\bfseries\fontsize{11}{13.2}\selectfont}
  {\thesubsubsection}{0.75em}{}
\titlespacing*{\subsubsection}{0pt}{0.5\baselineskip}{0.25\baselineskip}

\makeatletter
\renewcommand*\l@section{\@dottedtocline{1}{1.5em}{3.0em}}
\makeatother
\setcounter{tocdepth}{2}
\setcounter{secnumdepth}{2}
\renewcommand{\thesection}{\arabic{section}}
\renewcommand{\thesubsection}{\thesection.\arabic{subsection}}

\fancypagestyle{tesis}{
  \fancyhf{}
  \fancyhead[R]{\normalfont\fontsize{11}{13.2}\selectfont\thepage}
  \renewcommand{\headrulewidth}{0pt}
  \renewcommand{\footrulewidth}{0pt}
}
\fancypagestyle{plain}{
  \fancyhf{}
  \fancyhead[R]{\normalfont\fontsize{11}{13.2}\selectfont\thepage}
  \renewcommand{\headrulewidth}{0pt}
  \renewcommand{\footrulewidth}{0pt}
}
\pagestyle{tesis}

\newcommand{\tituloTesis}{Diseño y evaluación de una red neuronal informada por física para estimar la temperatura y la ampacidad en cables enterrados con heterogeneidad térmica}
\newcommand{\tituloInforme}{Informe de datos de casos de estudio y referencias de verificación térmica}

\begin{document}
\tesisbodyfont
\pagenumbering{arabic}

\begin{center}
{\fontsize{9}{10.8}\selectfont\MakeUppercase{\tituloTesis}\par}
\vspace{0.35cm}
{\fontsize{14}{16.8}\selectfont\bfseries\MakeUppercase{\tituloInforme}\par}
\vspace{0.35cm}
{\fontsize{11}{13.2}\selectfont\bfseries ELABORADO POR:\par}
\vspace{0.12cm}
{\fontsize{11}{13.2}\selectfont\bfseries LUIS ENRIQUE KOC GÓNGORA\par}
{\fontsize{11}{13.2}\selectfont\bfseries HERBERT ANTONIO MELÉNDEZ GARCÍA\par}
\end{center}

\vspace{0.45cm}
\noindent\textbf{Documento base:} Plan de tesis sobre el sistema cable--instalación--entorno térmico en cables eléctricos subterráneos.

\noindent\textbf{Propósito del informe:} organizar las instancias físicas del objeto de estudio y los casos normativos, numéricos, analíticos y manufacturados que verifican el artefacto PINN.


El informe organiza casos obtenidos de los PDF usados en el plan de
tesis, la matriz bibliográfica local y los documentos almacenados en
Zotero. La selección distingue entre instancias físicas del sistema
cable--instalación--entorno térmico y casos matemáticos o normativos
para verificar el artefacto PINN antes de aplicarlo al objeto físico.

El conjunto incluye fuentes con datos suficientes para reproducir una
instalación, un cálculo normativo o un problema de conducción de calor.
Para cada caso se priorizan la geometría, los materiales, las
propiedades térmicas, las condiciones de frontera y una solución de
referencia analítica, numérica o publicada.

\section{Resumen de fuentes con instancias
directas}\label{resumen-de-fuentes-con-instancias-directas}

{\def\LTcaptype{none} % do not increment counter
\begin{longtable}[]{@{}
  >{\RaggedRight\arraybackslash}p{0.13\linewidth}
  >{\RaggedRight\arraybackslash}p{0.14\linewidth}
  >{\RaggedRight\arraybackslash}p{0.34\linewidth}
  >{\RaggedRight\arraybackslash}p{0.33\linewidth}@{}}
\toprule\noalign{}
\begin{minipage}[b]{\linewidth}\raggedright
Clave
\end{minipage} & \begin{minipage}[b]{\linewidth}\raggedright
Fuente
\end{minipage} & \begin{minipage}[b]{\linewidth}\raggedright
Caso de instalación
\end{minipage} & \begin{minipage}[b]{\linewidth}\raggedright
Datos aprovechables
\end{minipage} \\
\midrule\noalign{}
\endhead
\bottomrule\noalign{}
\endlastfoot
kim\-2025 & Kim et al.~(2025) & UPCS 154 kV con seis cables, dos capas
planas, duct bank, bedding PAC/NC/arena y tres estratos de suelo
natural. & Cable multicapa, geometria de instalación, suelos, bedding,
fronteras térmicas, temperatura máxima y capacidad. \\
khumalo\-2025 & Khumalo et al.~(2025) & Cable MV XLPE bajo secado de
suelo; compara humedad, resistividad, profundidad, temperatura de suelo
y rating. & Resistividad del suelo, humedad, profundidad, temperatura
ambiente del suelo, ampacidad. \\
aldulaimi\-2024 & Al-Dulaimi et al.~(2024) & Cable 132 kV XLPE con
FEM-BPNN; arreglos plano/trébol, suelos mono/multicapa, backfill
SCMB/FTB y zonas secas. & Geometria 2D, configuraciones, materiales,
conductividades, condiciones de frontera, Tmax. \\
atoccsa\-2024 & Atoccsa et al.~(2024) & Cable 220 kV XLPE con backfill
térmico optimizado por PSO dinámico. & Datos del cable, variables de
zanja/backfill, costos, ampacidad con y sin backfill. \\
oclon\-2015 & Ocłoń et al.~(2015) & Sistema 400 kV en formación plana con
tubos HDPE, SBM y FTB optimizado por PSO. & Capas del cable, ductos,
conductividades, dominio FEM, profundidad, separación, Tmax. \\
aras\-2005 & Aras et al.~(2005) & Cable 154 kV XLPE comparado por IEC, FEM
y ensayo térmico de laboratorio. & Geometria, propiedades térmicas,
profundidad, dominio FEM, ampacidad/flujo térmico. \\
quan2019 & Quan et al.~(2019) & Tres cables enterrados en formación
plana o trébol, con suelo nativo o FTB y conductividad dependiente de T.
& Dominio 2D, geometría, fronteras, profundidad, separación, temperatura
y ampacidad. \\
oladunjoye2012 & Oladunjoye et al.~(2012) & Diez mediciones in situ de
suelo en una central eléctrica de Nigeria. & Resistividad,
conductividad, temperatura, humedad, densidad y porosidad del suelo. \\
mobius2025 & Möbius et al.~(2025) & Tres cables MV en trébol, enterrados
en Hamburgo y evaluados con 32 años de datos de suelo. & Composición y
humedad del suelo, temperatura, resistividad y ampacidad estacional. \\
cigre\-2022\_case0 & CIGRE TB 880 (2022) & Cable 132 kV enterrado; caso
base y 14 variantes de instalación y pérdidas. & Geometría, parámetros
IEC, pérdidas, temperaturas y ampacidades convergidas. \\
cigre\-2025\_case1 & CIGRE TB 963 (2025) & Tres cables 132 kV en ductos
HDPE, backfill y suelo; nueve variantes FEM 2D. & Dominio, cable
multicapa, materiales, fronteras, pérdidas y ampacidades. \\
iec\-60853\_annexA & IEC 60853-3 (2002) & Cable tripolar 132 kV sometido a
carga cíclica y secado parcial del suelo. & Dimensiones, resistencias y
capacitancias térmicas, ciclo de carga y corriente pico. \\
pan\-2025\_plate & Pan et al.~(2025) & Placa rectangular 2D sin fuente con
una condición de borde desconocida. & Dominio, PDE de Laplace,
fronteras, referencia FEM, muestreo y errores PINN. \\
cigre\-2025\_annulus & CIGRE TB 963 (2025) & Anillo conductor homogéneo
con flujo interior y temperatura exterior prescrita. & Solución
analítica completa para verificar el campo de temperatura. \\
hahn\-2012\_rect & Hahn y Özişik (2012) & Rectángulo 2D estacionario con
tres bordes a temperatura fija y uno convectivo. & Solución analítica en
serie para verificar una condición de Robin. \\
mms\-2d\_suite & Elaboración propia & Dos problemas manufacturados 2D:
conductividad variable e interfaz multimaterial. & Solución exacta,
fuente compatible, fronteras e interfaces para pruebas unitarias de la
PINN. \\
\end{longtable}
}

\section{Kim et al.~(2025) - UPCS 154 kV con bedding
PAC/NC/arena}\label{kim-et-al.-2025---upcs-154-kv-con-bedding-pacncarena}

\textbf{Referencia.} Kim, Y.-S., Cong, H. N., Dinh, B. H., \& Kim, H.-K.
(2025). Effect of ambient air and ground temperatures on heat transfer
in underground power cable system buried in newly developed cable
bedding material. \emph{Geothermics, 125}, 103151.
https://doi.org/10.1016/j.geothermics.2024.103151

Kim et al.~representan un sistema de cables multicapa enterrados en un
banco de ductos, rodeados por material de \textit{bedding} y varios
estratos de suelo. El caso relaciona la geometría y las propiedades
térmicas de la instalación con el clima local, la temperatura máxima y
la ampacidad.

{\def\LTcaptype{none} % do not increment counter
\begin{longtable}[]{@{}
  >{\raggedright\arraybackslash}p{(\linewidth - 4\tabcolsep) * \real{0.1060}}
  >{\raggedright\arraybackslash}p{(\linewidth - 4\tabcolsep) * \real{0.2781}}
  >{\raggedright\arraybackslash}p{(\linewidth - 4\tabcolsep) * \real{0.6159}}@{}}
\toprule\noalign{}
\begin{minipage}[b]{\linewidth}\raggedright
Bloque
\end{minipage} & \begin{minipage}[b]{\linewidth}\raggedright
Dato extraído
\end{minipage} & \begin{minipage}[b]{\linewidth}\raggedright
Valor o descripción
\end{minipage} \\
\midrule\noalign{}
\endhead
\bottomrule\noalign{}
\endlastfoot
Sistema & Tensión del cable & 154 kV \\
Sistema & Arreglo & Seis cables en dos capas de formación plana dentro
de duct bank. \\
Dominio & Entorno & Bloque de bedding + tres capas de suelo natural. \\
Frontera & Borde inferior & Temperatura constante del terreno: 15.2
°C. \\
Frontera & Superficie & Convección con aire ambiente; velocidad de
viento usada: 1.32 m/s invierno y 1.17 m/s verano. \\
Clima & Casos ambientales & Verano: aire 27.2 °C; invierno: aire 5.4
°C. \\
Criterio térmico & Temperatura máxima admisible del conductor & 90
°C. \\
Operación & Caso crítico reportado & Julio: 1026 A, aire 27.2 °C y
temperatura del suelo alrededor del cable 17.0 °C. \\
\end{longtable}
}

{\def\LTcaptype{none} % do not increment counter
\begin{longtable}[]{@{}
  >{\raggedright\arraybackslash}p{(\linewidth - 6\tabcolsep) * \real{0.2017}}
  >{\raggedright\arraybackslash}p{(\linewidth - 6\tabcolsep) * \real{0.2773}}
  >{\raggedright\arraybackslash}p{(\linewidth - 6\tabcolsep) * \real{0.3193}}
  >{\raggedright\arraybackslash}p{(\linewidth - 6\tabcolsep) * \real{0.2017}}@{}}
\toprule\noalign{}
\begin{minipage}[b]{\linewidth}\raggedright
Capa / dato del cable
\end{minipage} & \begin{minipage}[b]{\linewidth}\raggedright
Material
\end{minipage} & \begin{minipage}[b]{\linewidth}\raggedright
Dimensión
\end{minipage} & \begin{minipage}[b]{\linewidth}\raggedright
Conductividad k, W/(m K)
\end{minipage} \\
\midrule\noalign{}
\endhead
\bottomrule\noalign{}
\endlastfoot
Conductor & Cobre & dc = 42.4 mm; sección nominal 1200 mm² & 400 \\
Pantalla del conductor & Semiconductivo & d = 46.4 mm & 0.2857 \\
Aislamiento & XLPE & d = 80.4 mm & 0.2857 \\
Pantalla del aislamiento & Semiconductivo & d = 83.0 mm & 0.2857 \\
Cinta semiconductiva & Tape & d = 85.0 mm & 0.167 \\
Vaina metálica & Aluminio & d = 90.0 mm & 237 \\
Cubierta externa & PE & d = 100.0 mm & 0.2857 \\
Tubo/casing & PE & d = 220 mm; espesor 10 mm & 0.2857 \\
Dato eléctrico & Resistencia DC a 20 °C & 0.0151 Ω/km & \\
Dato eléctrico & Coeficiente térmico del conductor & 0.0393 & \\
\end{longtable}
}

{\def\LTcaptype{none} % do not increment counter
\begin{longtable}[]{@{}
  >{\raggedright\arraybackslash}p{(\linewidth - 10\tabcolsep) * \real{0.0690}}
  >{\raggedright\arraybackslash}p{(\linewidth - 10\tabcolsep) * \real{0.0460}}
  >{\raggedright\arraybackslash}p{(\linewidth - 10\tabcolsep) * \real{0.3218}}
  >{\raggedright\arraybackslash}p{(\linewidth - 10\tabcolsep) * \real{0.1149}}
  >{\raggedright\arraybackslash}p{(\linewidth - 10\tabcolsep) * \real{0.1149}}
  >{\raggedright\arraybackslash}p{(\linewidth - 10\tabcolsep) * \real{0.3333}}@{}}
\toprule\noalign{}
\begin{minipage}[b]{\linewidth}\raggedright
Suelo
\end{minipage} & \begin{minipage}[b]{\linewidth}\raggedright
USCS
\end{minipage} & \begin{minipage}[b]{\linewidth}\raggedright
Peso unitario natural, kN/m³
\end{minipage} & \begin{minipage}[b]{\linewidth}\raggedright
Humedad, \%
\end{minipage} & \begin{minipage}[b]{\linewidth}\raggedright
k, W/(m K)
\end{minipage} & \begin{minipage}[b]{\linewidth}\raggedright
Resistividad térmica, °C cm/W
\end{minipage} \\
\midrule\noalign{}
\endhead
\bottomrule\noalign{}
\endlastfoot
Capa 1 & SC & 18.081 & 23.25 & 1.804 & 55.44 \\
Capa 2 & CL & 18.884 & 26.27 & 1.351 & 74.01 \\
Capa 3 & CL & 19.987 & 23.14 & 1.517 & 65.94 \\
\end{longtable}
}

{\def\LTcaptype{none} % do not increment counter
\begin{longtable}[]{@{}
  >{\raggedright\arraybackslash}p{(\linewidth - 8\tabcolsep) * \real{0.2456}}
  >{\raggedright\arraybackslash}p{(\linewidth - 8\tabcolsep) * \real{0.1754}}
  >{\raggedright\arraybackslash}p{(\linewidth - 8\tabcolsep) * \real{0.2544}}
  >{\raggedright\arraybackslash}p{(\linewidth - 8\tabcolsep) * \real{0.2368}}
  >{\raggedright\arraybackslash}p{(\linewidth - 8\tabcolsep) * \real{0.0877}}@{}}
\toprule\noalign{}
\begin{minipage}[b]{\linewidth}\raggedright
Material de bedding/backfill
\end{minipage} & \begin{minipage}[b]{\linewidth}\raggedright
Densidad seca, kg/m³
\end{minipage} & \begin{minipage}[b]{\linewidth}\raggedright
Fluidez / bleeding
\end{minipage} & \begin{minipage}[b]{\linewidth}\raggedright
Resistencia compresiva, MPa
\end{minipage} & \begin{minipage}[b]{\linewidth}\raggedright
k, W/(m K)
\end{minipage} \\
\midrule\noalign{}
\endhead
\bottomrule\noalign{}
\endlastfoot
Arena natural & 1603 & No reportada & No reportada & 1.365 \\
Concreto normal (NC) & 2093 & Fluidez 180 mm & 15.0 & 2.093 \\
PAC & 1410 & Bleeding 1.7\%; grout 382 mm & 1.56 & 2.094 \\
CLSM-RS10 & 2735 & Bleeding 3.0\%; fluidez 235 mm & 2.735 & 2.150 \\
\end{longtable}
}

{\def\LTcaptype{none} % do not increment counter
\begin{longtable}[]{@{}
  >{\raggedright\arraybackslash}p{(\linewidth - 6\tabcolsep) * \real{0.3712}}
  >{\raggedright\arraybackslash}p{(\linewidth - 6\tabcolsep) * \real{0.0985}}
  >{\raggedright\arraybackslash}p{(\linewidth - 6\tabcolsep) * \real{0.1742}}
  >{\raggedright\arraybackslash}p{(\linewidth - 6\tabcolsep) * \real{0.3561}}@{}}
\toprule\noalign{}
\begin{minipage}[b]{\linewidth}\raggedright
Condición
\end{minipage} & \begin{minipage}[b]{\linewidth}\raggedright
Arena natural
\end{minipage} & \begin{minipage}[b]{\linewidth}\raggedright
PAC / NC
\end{minipage} & \begin{minipage}[b]{\linewidth}\raggedright
Dato útil para la tesis
\end{minipage} \\
\midrule\noalign{}
\endhead
\bottomrule\noalign{}
\endlastfoot
Temperatura máxima del conductor en verano & 77.6 °C & 70.6 °C &
Backfill de mayor k reduce Tmax en 7.0 °C. \\
Temperatura máxima del conductor en invierno & 57.1 °C & 49.9 °C &
Reducción aproximada de 7.2 °C. \\
Diferencia conductor central-superior en verano & 1.0 °C & 1.4 °C &
Evidencia de gradiente interno por posición. \\
Diferencia conductor central-superior en invierno & 1.3 °C & 1.6 °C &
Confirma sensibilidad geométrica del arreglo. \\
Capacidad en caso crítico & No aplica & 1026 A con Tmax 70.6 °C &
Escenario calibrable para comparación FEM/PINN. \\
\end{longtable}
}

\textbf{Aporte a la tesis.} La investigación usa este caso para evaluar
rellenos con distinta conductividad y condiciones climáticas en la
frontera superior. La combinación de un cable XLPE multicapa, estratos
de suelo y resultados térmicos publicados permite comparar la respuesta
del artefacto PINN con una instalación físicamente definida.

\section{Khumalo et al.~(2025) - rating bajo secado de
suelo}\label{khumalo-et-al.-2025---rating-bajo-secado-de-suelo}

\textbf{Referencia.} Khumalo, N. Q., Naidoo, R. M., Mbungu, N. T., \&
Bansal, R. C. (2025). A critical assessment of cable rating methods
under soil drying out conditions. \emph{International Transactions on
Electrical Energy Systems, 2025}, Article 5946564.
https://doi.org/10.1155/etep/5946564

Khumalo et al.~cuantifican la relación entre el estado hídrico del
suelo, la resistividad térmica y la ampacidad. El estudio también
analiza la sensibilidad de la respuesta frente a la profundidad de
instalación y la temperatura del terreno.

{\def\LTcaptype{none} % do not increment counter
\begin{longtable}[]{@{}
  >{\raggedright\arraybackslash}p{(\linewidth - 4\tabcolsep) * \real{0.0975}}
  >{\raggedright\arraybackslash}p{(\linewidth - 4\tabcolsep) * \real{0.1271}}
  >{\raggedright\arraybackslash}p{(\linewidth - 4\tabcolsep) * \real{0.7754}}@{}}
\toprule\noalign{}
\begin{minipage}[b]{\linewidth}\raggedright
Componente
\end{minipage} & \begin{minipage}[b]{\linewidth}\raggedright
Dato extraído
\end{minipage} & \begin{minipage}[b]{\linewidth}\raggedright
Valor o descripción
\end{minipage} \\
\midrule\noalign{}
\endhead
\bottomrule\noalign{}
\endlastfoot
Cable & Tipo & Cable MV XLPE de tres núcleos. \\
Cable & Capas identificadas & Conductor de cobre, pantalla
semiconductiva, XLPE, pantalla del núcleo, cinta semiconductiva,
pantalla de cobre, relleno, bedding FR-PVC, armadura de acero
galvanizado y cubierta PVC. \\
Condición de referencia & Resistividad térmica del suelo & 1.2 K m/W. \\
Condición de referencia & Temperatura ambiente del suelo & 25 °C. \\
Condición de referencia & Profundidad de tendido & 850 mm. \\
\end{longtable}
}

{\def\LTcaptype{none} % do not increment counter
\begin{longtable}[]{@{}
  >{\raggedright\arraybackslash}p{(\linewidth - 4\tabcolsep) * \real{0.2439}}
  >{\raggedright\arraybackslash}p{(\linewidth - 4\tabcolsep) * \real{0.3293}}
  >{\raggedright\arraybackslash}p{(\linewidth - 4\tabcolsep) * \real{0.4268}}@{}}
\toprule\noalign{}
\begin{minipage}[b]{\linewidth}\raggedright
Humedad del suelo, \%
\end{minipage} & \begin{minipage}[b]{\linewidth}\raggedright
Resistividad térmica, K m/W
\end{minipage} & \begin{minipage}[b]{\linewidth}\raggedright
Lectura para el objeto de estudio
\end{minipage} \\
\midrule\noalign{}
\endhead
\bottomrule\noalign{}
\endlastfoot
14.5 & 0.596 & Suelo húmedo / condición favorable. \\
5 & 1.646 & Secado parcial relevante. \\
2 & 2.341 & Secado severo. \\
0 & 3.720 & Suelo seco / condición crítica. \\
\end{longtable}
}

{\def\LTcaptype{none} % do not increment counter
\begin{longtable}[]{@{}
  >{\raggedright\arraybackslash}p{(\linewidth - 10\tabcolsep) * \real{0.3182}}
  >{\raggedright\arraybackslash}p{(\linewidth - 10\tabcolsep) * \real{0.2159}}
  >{\raggedright\arraybackslash}p{(\linewidth - 10\tabcolsep) * \real{0.1136}}
  >{\raggedright\arraybackslash}p{(\linewidth - 10\tabcolsep) * \real{0.1136}}
  >{\raggedright\arraybackslash}p{(\linewidth - 10\tabcolsep) * \real{0.1136}}
  >{\raggedright\arraybackslash}p{(\linewidth - 10\tabcolsep) * \real{0.1250}}@{}}
\toprule\noalign{}
\begin{minipage}[b]{\linewidth}\raggedright
Condición
\end{minipage} & \begin{minipage}[b]{\linewidth}\raggedright
Resistividad, K m/W
\end{minipage} & \begin{minipage}[b]{\linewidth}\raggedright
Cable A, A
\end{minipage} & \begin{minipage}[b]{\linewidth}\raggedright
Cable B, A
\end{minipage} & \begin{minipage}[b]{\linewidth}\raggedright
Cable C, A
\end{minipage} & \begin{minipage}[b]{\linewidth}\raggedright
Promedio, A
\end{minipage} \\
\midrule\noalign{}
\endhead
\bottomrule\noalign{}
\endlastfoot
Suelo 14.5\% humedad & 0.596 & 518.746 & 519.540 & 516.723 & 518.34 \\
Suelo importado / referencia & 1.200 & 382.371 & 382.869 & 381.143 &
382.13 \\
Suelo 5\% humedad & 1.646 & 330.663 & 331.070 & 329.667 & 330.47 \\
Suelo 2\% humedad & 2.341 & 280.169 & 280.495 & 279.372 & 280.01 \\
Suelo seco & 3.720 & 224.330 & 224.582 & 223.730 & 224.21 \\
\end{longtable}
}

{\def\LTcaptype{none} % do not increment counter
\begin{longtable}[]{@{}
  >{\raggedright\arraybackslash}p{(\linewidth - 4\tabcolsep) * \real{0.2500}}
  >{\raggedright\arraybackslash}p{(\linewidth - 4\tabcolsep) * \real{0.2417}}
  >{\raggedright\arraybackslash}p{(\linewidth - 4\tabcolsep) * \real{0.5083}}@{}}
\toprule\noalign{}
\begin{minipage}[b]{\linewidth}\raggedright
Variable
\end{minipage} & \begin{minipage}[b]{\linewidth}\raggedright
Valores evaluados
\end{minipage} & \begin{minipage}[b]{\linewidth}\raggedright
Resultado extraído
\end{minipage} \\
\midrule\noalign{}
\endhead
\bottomrule\noalign{}
\endlastfoot
Temperatura ambiente del suelo & 24, 25, 26, 27 y 28 °C & Para Cable A:
386.597, 382.371, 378.099, 373.777 y 369.405 A. \\
Profundidad de tendido & 750, 850, 950, 1000 y 1150 mm & Para Cable A:
382.321, 382.371, 377.308, 375.040 y 369.062 A. \\
Peor caso & 3.72 K m/W; 1150 mm; 28 °C & Cable A/B/C: 208.70, 208.93 y
208.19 A. \\
Caso ideal & 1.2 K m/W; 850 mm; 25 °C & Cable A/B/C: 382.37, 382.87 y
381.14 A. \\
Reducción promedio & Peor caso frente a referencia & 45.40\%. \\
\end{longtable}
}

\textbf{Aporte a la tesis.} Estos datos permiten construir escenarios de
baja conductividad y alta resistividad asociados con el secado. Así, el
artefacto puede evaluarse frente a contrastes de \(k(x,y)\) y cambios de
ampacidad que una representación homogénea fija no describe.

\section{Al-Dulaimi et al.~(2024) - FEM-BPNN con suelos y backfill
variables}\label{al-dulaimi-et-al.-2024---fem-bpnn-con-suelos-y-backfill-variables}

\textbf{Referencia.} Al-Dulaimi, A. A., Guneser, M. T., Hameed, A. A.,
García Márquez, F. P., \& Gouda, O. E. (2024). Adaptive FEM-BPNN model
for predicting underground cable temperature considering varied soil
composition. \emph{Engineering Science and Technology, an International
Journal, 51}, 101658. https://doi.org/10.1016/j.jestch.2024.101658

Al-Dulaimi et al.~convierten varios escenarios FEM de cables enterrados
en una base de aprendizaje para una red neuronal. Aunque la red no es
una PINN, sus entradas y salidas reúnen las variables del estudio:
geometría, configuración, profundidad, separación, \textit{backfill},
suelo y temperatura máxima.

{\def\LTcaptype{none} % do not increment counter
\begin{longtable}[]{@{}
  >{\raggedright\arraybackslash}p{(\linewidth - 4\tabcolsep) * \real{0.2336}}
  >{\raggedright\arraybackslash}p{(\linewidth - 4\tabcolsep) * \real{0.2991}}
  >{\raggedright\arraybackslash}p{(\linewidth - 4\tabcolsep) * \real{0.4673}}@{}}
\toprule\noalign{}
\begin{minipage}[b]{\linewidth}\raggedright
Bloque
\end{minipage} & \begin{minipage}[b]{\linewidth}\raggedright
Dato extraído
\end{minipage} & \begin{minipage}[b]{\linewidth}\raggedright
Valor o descripción
\end{minipage} \\
\midrule\noalign{}
\endhead
\bottomrule\noalign{}
\endlastfoot
Cable & Tensión & 132 kV. \\
Cable & Sección geométrica del conductor & 749.9 mm². \\
Criterio térmico & Temperatura máxima admisible & 90 °C para conductor
XLPE en régimen estacionario. \\
Dominio FEM & Área de simulación & 4 m x 4 m. \\
Backfill & Ancho del bloque & 2 m. \\
Profundidad H & Valores & 0.8, 1.0 y 1.2 m. \\
Separación S & Formación plana & 0.1 a 0.4 m. \\
Separación S & Formación trébol & 0.008 a 0.2 m según el artículo. \\
Temperatura superficial & Tsg & 20 a 50 °C. \\
Temperatura inicial/suelo & Tsoil & 14 °C. \\
Carga & Máxima reportada para dataset & 945 A. \\
\end{longtable}
}

{\def\LTcaptype{none} % do not increment counter
\begin{longtable}[]{@{}
  >{\raggedright\arraybackslash}p{(\linewidth - 4\tabcolsep) * \real{0.2710}}
  >{\raggedright\arraybackslash}p{(\linewidth - 4\tabcolsep) * \real{0.2243}}
  >{\raggedright\arraybackslash}p{(\linewidth - 4\tabcolsep) * \real{0.5047}}@{}}
\toprule\noalign{}
\begin{minipage}[b]{\linewidth}\raggedright
Material / zona
\end{minipage} & \begin{minipage}[b]{\linewidth}\raggedright
Conductividad k, W/(m K)
\end{minipage} & \begin{minipage}[b]{\linewidth}\raggedright
Observación
\end{minipage} \\
\midrule\noalign{}
\endhead
\bottomrule\noalign{}
\endlastfoot
Conductor de cobre & 386 & Capa interna conductora. \\
Aislamiento XLPE & 0.2875 & Capa dieléctrica. \\
Cubierta HDPE & 0.2875 & Cubierta externa. \\
Pantalla/vaina de cobre & 386 & Ruta metálica de alta k. \\
Suelo medio & 2.28 & Condición de suelo usada como referencia. \\
SCMB & 1.00 & Sand-cement mixture backfill. \\
FTB & 1.54 & Fluidized thermal backfill. \\
Suelo nativo / sand clay loam & Variable & El estudio evalúa estados
húmedos/secos y composiciones. \\
\end{longtable}
}

{\def\LTcaptype{none} % do not increment counter
\begin{longtable}[]{@{}
  >{\raggedright\arraybackslash}p{(\linewidth - 6\tabcolsep) * \real{0.3271}}
  >{\raggedright\arraybackslash}p{(\linewidth - 6\tabcolsep) * \real{0.0748}}
  >{\raggedright\arraybackslash}p{(\linewidth - 6\tabcolsep) * \real{0.1682}}
  >{\raggedright\arraybackslash}p{(\linewidth - 6\tabcolsep) * \real{0.4299}}@{}}
\toprule\noalign{}
\begin{minipage}[b]{\linewidth}\raggedright
Escenario FEM
\end{minipage} & \begin{minipage}[b]{\linewidth}\raggedright
Backfill
\end{minipage} & \begin{minipage}[b]{\linewidth}\raggedright
Tmax reportada, °C
\end{minipage} & \begin{minipage}[b]{\linewidth}\raggedright
Lectura
\end{minipage} \\
\midrule\noalign{}
\endhead
\bottomrule\noalign{}
\endlastfoot
Una capa homogénea, formación plana & SCMB & 84.31 & Por debajo del
límite de 90 °C. \\
Una capa homogénea, formación plana & FTB & 82.99 & FTB reduce Tmax
frente a SCMB. \\
Una capa homogénea, trébol & SCMB & 86.10 & Mayor Tmax por proximidad
térmica. \\
Una capa homogénea, trébol & FTB & 84.64 & Reducción con FTB, pero sigue
mayor que plana. \\
Multicapa homogénea, plana & SCMB & 82.58 & La estratificación cambia el
resultado. \\
Multicapa homogénea, plana & FTB & 81.22 & Mejor disipación. \\
Multicapa homogénea, trébol & SCMB & 84.64 & Mayor concentración
térmica. \\
Multicapa homogénea, trébol & FTB & 83.15 & Mejor que SCMB. \\
Multicapa no homogénea, plana & SCMB & 83.65 & Caso heterogéneo
explícito. \\
Multicapa no homogénea, plana & FTB & 82.42 & Backfill de mayor k reduce
Tmax. \\
Multicapa no homogénea, trébol & SCMB & 86.26 & Caso más exigente entre
los reportados. \\
Multicapa no homogénea, trébol & FTB & 84.89 & FTB reduce el máximo. \\
\end{longtable}
}

{\def\LTcaptype{none} % do not increment counter
\begin{longtable}[]{@{}
  >{\raggedright\arraybackslash}p{(\linewidth - 2\tabcolsep) * \real{0.2381}}
  >{\raggedright\arraybackslash}p{(\linewidth - 2\tabcolsep) * \real{0.7619}}@{}}
\toprule\noalign{}
\begin{minipage}[b]{\linewidth}\raggedright
Variable del conjunto de datos
\end{minipage} & \begin{minipage}[b]{\linewidth}\raggedright
Uso en el artefacto PINN o en el caso de verificación
\end{minipage} \\
\midrule\noalign{}
\endhead
\bottomrule\noalign{}
\endlastfoot
Tsg & Condición de frontera térmica superior. \\
S & Distancia entre cables; controla proximidad de fuentes. \\
H & Profundidad de instalación. \\
kbackfill & Contraste térmico entre relleno y suelo. \\
Código de entorno & Identifica suelo homogéneo, multicapa, no homogéneo
o zona seca. \\
Tmax & Salida escalar para comparar con campo térmico 2D. \\
\end{longtable}
}

\textbf{Aporte a la tesis.} La estructura paramétrica del caso permite
variar la geometría, el patrón de heterogeneidad y el material de
relleno. La investigación puede usar esas variaciones para comparar la
temperatura máxima estimada por el artefacto PINN con una referencia
FEM.

\section{Atoccsa et al.~(2024) - optimización de ampacidad con backfill
térmico}\label{atoccsa-et-al.-2024---optimizaciuxf3n-de-ampacidad-con-backfill-tuxe9rmico}

\textbf{Referencia.} Atoccsa, B. A., Puma, D. W., Mendoza, D., Urday,
E., Ronceros, C., \& Palma, M. T. (2024). Optimization of ampacity in
high-voltage underground cables with thermal backfill using dynamic PSO
and adaptive strategies. \emph{Energies, 17}(5), 1023.
https://doi.org/10.3390/en17051023

Atoccsa et al.~analizan un caso de diseño en el que el \textit{backfill}
térmico se optimiza junto con la geometría de instalación. Los
resultados muestran cómo la modificación del entorno próximo al cable
cambia la ampacidad de un sistema de alta tensión.

{\def\LTcaptype{none} % do not increment counter
\begin{longtable}[]{@{}
  >{\raggedright\arraybackslash}p{(\linewidth - 2\tabcolsep) * \real{0.4667}}
  >{\raggedright\arraybackslash}p{(\linewidth - 2\tabcolsep) * \real{0.5333}}@{}}
\toprule\noalign{}
\begin{minipage}[b]{\linewidth}\raggedright
Dato del cable 220 kV XLPE
\end{minipage} & \begin{minipage}[b]{\linewidth}\raggedright
Valor
\end{minipage} \\
\midrule\noalign{}
\endhead
\bottomrule\noalign{}
\endlastfoot
Conductor & Cobre Milliken de 5 segmentos, recocido. \\
Sección del conductor & 2000 mm². \\
Diámetro del conductor dc & 54.5 mm. \\
Espesor de pantalla semiconductiva & 3.5 mm. \\
Espesor de aislamiento XLPE & 24.0 mm. \\
Diámetro externo del aislamiento Di & 107.1 mm. \\
Espesor de vaina de aluminio & 2.8 mm. \\
Diámetro externo de vaina Ds & 137.4 mm. \\
Espesor de cubierta externa & 5.0 mm. \\
Diámetro externo del cable De & 147.7 mm. \\
Temperatura máxima del conductor & 90 °C. \\
Frecuencia & 60 Hz. \\
Resistencia del conductor a 20 °C & 0.009 Ω/km. \\
Tensión nominal fase-fase & 220 kV. \\
\end{longtable}
}

{\def\LTcaptype{none} % do not increment counter
\begin{longtable}[]{@{}
  >{\raggedright\arraybackslash}p{(\linewidth - 6\tabcolsep) * \real{0.2169}}
  >{\raggedright\arraybackslash}p{(\linewidth - 6\tabcolsep) * \real{0.4217}}
  >{\raggedright\arraybackslash}p{(\linewidth - 6\tabcolsep) * \real{0.1807}}
  >{\raggedright\arraybackslash}p{(\linewidth - 6\tabcolsep) * \real{0.1807}}@{}}
\toprule\noalign{}
\begin{minipage}[b]{\linewidth}\raggedright
Variable de diseño
\end{minipage} & \begin{minipage}[b]{\linewidth}\raggedright
Significado
\end{minipage} & \begin{minipage}[b]{\linewidth}\raggedright
Límite inferior
\end{minipage} & \begin{minipage}[b]{\linewidth}\raggedright
Límite superior
\end{minipage} \\
\midrule\noalign{}
\endhead
\bottomrule\noalign{}
\endlastfoot
L & Profundidad del cable & 0.5 m & 2.0 m \\
LG & Profundidad del centro del backfill & 0.6 m & 4.0 m \\
w & Ancho del backfill & 1.2 m & 4.0 m \\
h & Espesor/altura del backfill & 0.6 m & 3.0 m \\
s & Separación entre cables & De ≈ 0.147 m & 2.0 m \\
s1 & Separación auxiliar & 0.3 m & 2.0 m \\
\end{longtable}
}

{\def\LTcaptype{none} % do not increment counter
\begin{longtable}[]{@{}
  >{\raggedright\arraybackslash}p{(\linewidth - 4\tabcolsep) * \real{0.3846}}
  >{\raggedright\arraybackslash}p{(\linewidth - 4\tabcolsep) * \real{0.2154}}
  >{\raggedright\arraybackslash}p{(\linewidth - 4\tabcolsep) * \real{0.4000}}@{}}
\toprule\noalign{}
\begin{minipage}[b]{\linewidth}\raggedright
Componente de costo
\end{minipage} & \begin{minipage}[b]{\linewidth}\raggedright
Costo unitario
\end{minipage} & \begin{minipage}[b]{\linewidth}\raggedright
Expresión geométrica usada
\end{minipage} \\
\midrule\noalign{}
\endhead
\bottomrule\noalign{}
\endlastfoot
Excavación & 16.5 USD/m³ & w·LG + w·h/2 \\
Remoción de tierra & 13.15 USD/m³ & w·LG + w·h/2 \\
Backfill de arena térmica & 28.5 USD/m³ & w·h - (3/4)πDe² \\
\end{longtable}
}

{\def\LTcaptype{none} % do not increment counter
\begin{longtable}[]{@{}ll@{}}
\toprule\noalign{}
Resultado optimizado & Valor \\
\midrule\noalign{}
\endhead
\bottomrule\noalign{}
\endlastfoot
L & 0.500 m \\
LG & 0.872 m \\
w & 3.562 m \\
h & 1.344 m \\
s & 1.481 m \\
λ1 & 2.667 \\
Costo total & 300 USD \\
Costo de backfill & 94.7 USD \\
Ampacidad con backfill & 1156.915 A \\
Ampacidad sin backfill & 969.9 A \\
Incremento aproximado & 18.45\% respecto al caso sin backfill. \\
Pérdidas dieléctricas Wd & 3 x 3.546 W/m \\
Pérdidas de carga Wl & 3 x 17.67 W/m \\
\end{longtable}
}

\textbf{Aporte a la tesis.} Este caso incorpora la ampacidad como
respuesta de una geometría y un \textit{backfill} definidos, además de
la temperatura máxima. Por tanto, permite verificar si el artefacto
reproduce las tendencias de diseño asociadas con la conductividad
térmica y las dimensiones del relleno.

\section{Ocłoń et al.~(2015) - bedding FTB optimizado en sistema 400
kV}\label{ocux142oux144-et-al.-2015---bedding-ftb-optimizado-en-sistema-400-kv}

\textbf{Referencia.} Ocłoń, P., Cisek, P., Taler, D., Pilarczyk, M., \&
Szwarc, T. (2015). Optimizing of the underground power cable bedding
using momentum-type particle swarm optimization method. \emph{Energy,
92}(2), 230--239. https://doi.org/10.1016/j.energy.2015.04.100

Ocłoń et al.~modelan tres cables subterráneos de 400 kV en formación
plana, instalados en tubos HDPE con mezcla SBM y rodeados por
\textit{backfill} FTB. La geometría separa el suelo nativo, el ducto, el
relleno interior y el \textit{bedding} térmico exterior.

{\def\LTcaptype{none} % do not increment counter
\begin{longtable}[]{@{}
  >{\raggedright\arraybackslash}p{(\linewidth - 2\tabcolsep) * \real{0.1739}}
  >{\raggedright\arraybackslash}p{(\linewidth - 2\tabcolsep) * \real{0.8261}}@{}}
\toprule\noalign{}
\begin{minipage}[b]{\linewidth}\raggedright
Elemento del sistema
\end{minipage} & \begin{minipage}[b]{\linewidth}\raggedright
Dato extraído
\end{minipage} \\
\midrule\noalign{}
\endhead
\bottomrule\noalign{}
\endlastfoot
Configuración & Tres cables subterráneos de 400 kV en formación
plana/in-line. \\
Ducto & Tubos HDPE con diámetro externo 278 mm y espesor 14 mm. \\
Relleno dentro del ducto & SBM: mezcla sand-bentonite, densidad 1700
kg/m³, k = 0.95 W/(m K). \\
Bedding exterior & FTB: mezcla SGFC con 41\% agregado fino, 49\%
agregado grueso, 2.5\% cemento y 7.5\% fly ash; densidad seca 2187
kg/m³; k = 1.54 W/(m K). \\
Suelo madre & k = 1.00 W/(m K). \\
Ducto HDPE & k = 0.48 W/(m K). \\
Dominio FEM & Modelo 2D estacionario; dominio cuadrado de 10 m x 10 m
usando simetría. \\
Profundidad & H = 2 m desde el nivel de cruce vial. \\
Frontera superior & Temperatura Tg = 30 °C. \\
Fronteras laterales/inferior & Adiabáticas por simetría/aislamiento
térmico en el dominio. \\
Carga máxima asumida & 1145 A. \\
Criterio de operación & Topt = 65 °C; el máximo del conductor no debe
exceder ese valor. \\
\end{longtable}
}

{\def\LTcaptype{none} % do not increment counter
\begin{longtable}[]{@{}ll@{}}
\toprule\noalign{}
Dato del cable 400 kV & Valor \\
\midrule\noalign{}
\endhead
\bottomrule\noalign{}
\endlastfoot
Sección del conductor & 1600 mm². \\
Diámetro del conductor & 49.6 mm. \\
Espesor total de aislamiento & 27 mm. \\
Diámetro externo & 127.9 mm. \\
Resistencia DC a 20 °C & 0.0113 Ω/km. \\
Resistencia AC a 65 °C y 50 Hz & 0.0157 Ω/km. \\
Current loading in ground in-line a 65 °C & 1145 A. \\
\end{longtable}
}

{\def\LTcaptype{none} % do not increment counter
\begin{longtable}[]{@{}llll@{}}
\toprule\noalign{}
Capa del cable & Material & Espesor / radio & k, W/(m K) \\
\midrule\noalign{}
\endhead
\bottomrule\noalign{}
\endlastfoot
Conductor & Cobre & dc = 49.6 mm; rc = 24.8 mm & 400 \\
Aislamiento & XLPE & 30.5 mm; rins = 55.3 mm & 0.2875 \\
Vaina/pantalla & Cobre & 6.4 mm; rsh = 61.7 mm & 400 \\
Cubierta & HDPE & 5.1 mm; rj = 66.8 mm & 0.2875 \\
\end{longtable}
}

{\def\LTcaptype{none} % do not increment counter
\begin{longtable}[]{@{}
  >{\raggedright\arraybackslash}p{(\linewidth - 2\tabcolsep) * \real{0.2754}}
  >{\raggedright\arraybackslash}p{(\linewidth - 2\tabcolsep) * \real{0.7246}}@{}}
\toprule\noalign{}
\begin{minipage}[b]{\linewidth}\raggedright
Variable optimizada
\end{minipage} & \begin{minipage}[b]{\linewidth}\raggedright
Rango / resultado
\end{minipage} \\
\midrule\noalign{}
\endhead
\bottomrule\noalign{}
\endlastfoot
l, separación entre ejes de cables & Rango 0.3 a 0.6 m; óptimo 0.600
m. \\
s, distancia al borde derecho del FTB & Rango 0.2 a 0.4 m; óptimo
0.2000-0.20044 m. \\
b, distancia al borde superior del FTB & Rango 0.2 a 0.4 m; óptimo
0.31327-0.31366 m. \\
p, distancia al borde inferior del FTB & Rango 0.2 a 0.4 m; óptimo
0.2000-0.20044 m. \\
Tmax & 65 °C en el conductor central. \\
Área media FTB modelada Ab & 0.3195-0.3196 m²; el área real por simetría
es aprox. 0.6390 m². \\
Interpretación del resultado & El conductor central queda más caliente
por interacción térmica y peor disipación que los laterales. \\
\end{longtable}
}

\textbf{Aporte a la tesis.} La separación de materiales permite estudiar
la proximidad entre fuentes térmicas y rellenos localizados. A partir de
esta geometría se construyen escenarios 2D con contrastes entre suelo
nativo, FTB, SBM y HDPE, de acuerdo con la heterogeneidad espacial
definida en el plan.

\section{Aras et al.~(2005) - comparación IEC/FEM/ensayo en cable 154
kV}\label{aras-et-al.-2005---comparaciuxf3n-iecfemensayo-en-cable-154-kv}

\textbf{Referencia.} Aras, F., Oysu, C., \& Yilmaz, G. (2005). An
assessment of the methods for calculating ampacity of underground power
cables. \emph{Electric Power Components and Systems, 33}(12),
1385--1402. https://doi.org/10.1080/15325000590964425

Aras et al.~comparan métodos de cálculo de ampacidad con FEM y con un
ensayo térmico. Aunque la instalación no incorpora el detalle de suelos
y rellenos de los estudios posteriores, sus condiciones homogéneas
proporcionan una referencia inicial para la verificación.

{\def\LTcaptype{none} % do not increment counter
\begin{longtable}[]{@{}
  >{\raggedright\arraybackslash}p{(\linewidth - 2\tabcolsep) * \real{0.3956}}
  >{\raggedright\arraybackslash}p{(\linewidth - 2\tabcolsep) * \real{0.6044}}@{}}
\toprule\noalign{}
\begin{minipage}[b]{\linewidth}\raggedright
Variable
\end{minipage} & \begin{minipage}[b]{\linewidth}\raggedright
Dato extraído
\end{minipage} \\
\midrule\noalign{}
\endhead
\bottomrule\noalign{}
\endlastfoot
Cable & 154 kV XLPE. \\
Temperatura de suelo & 20 °C para el caso del norte de Turquía. \\
Profundidad de enterramiento & 1.2 m. \\
Medio & Suelo homogéneo. \\
Temperatura máxima de operación XLPE & 90 °C. \\
Dominio FEM & 18 m de ancho x 10 m de profundidad; fronteras a 20 °C. \\
Pérdida dieléctrica & 3.57 W/m. \\
Conductividad del suelo & 1.2 W/(m K). \\
Conductividad XLPE & 0.2857 W/(m K). \\
Conductividad pantalla & 384.6 W/(m K). \\
\end{longtable}
}

{\def\LTcaptype{none} % do not increment counter
\begin{longtable}[]{@{}ll@{}}
\toprule\noalign{}
Dimensión del cable & Valor \\
\midrule\noalign{}
\endhead
\bottomrule\noalign{}
\endlastfoot
Diámetro del conductor & 37.7 mm. \\
Diámetro XLPE + semiconductor & 81.7 mm. \\
Diámetro de pantalla & 98.7 mm. \\
Diámetro de cubierta & 106.7 mm. \\
Profundidad h & 1200 mm. \\
\end{longtable}
}

{\def\LTcaptype{none} % do not increment counter
\begin{longtable}[]{@{}
  >{\raggedright\arraybackslash}p{(\linewidth - 2\tabcolsep) * \real{0.1707}}
  >{\raggedright\arraybackslash}p{(\linewidth - 2\tabcolsep) * \real{0.8293}}@{}}
\toprule\noalign{}
\begin{minipage}[b]{\linewidth}\raggedright
Resultado / referencia
\end{minipage} & \begin{minipage}[b]{\linewidth}\raggedright
Dato
\end{minipage} \\
\midrule\noalign{}
\endhead
\bottomrule\noalign{}
\endlastfoot
Mallas FEM ensayadas & 557, 690, 1074 y 1486 elementos; se selecciona
1074 elementos. \\
Comparación FEM vs IEC & Diferencia aproximada del orden de 1\% para el
caso principal. \\
Sensibilidad del aislamiento & Reducir el espesor de XLPE de 22 mm a 17
mm incrementa la ampacidad en 2.9\%. \\
Flujo térmico de referencia & Para 22 mm de aislamiento XLPE se reporta
0.585 W/m² en la superficie del conductor. \\
Ensayo experimental & 15 m de cable 154 kV XLPE; ciclos térmicos
aproximados de 310 K a 380 K; suelo representado por una capa de papel
con k = 0.475 W/(m K). \\
\end{longtable}
}

\textbf{Aporte a la tesis.} La investigación usa esta instalación como
caso base con suelo homogéneo y geometría simple. Primero se verifica la
conducción radial y enterrada; luego se incorporan \textit{backfill},
estratos o zonas secas para evaluar la respuesta frente a una
complejidad creciente.

\section{Quan et al.~(2019) - formaciones plana y trébol con suelo o
FTB}\label{quan-et-al.-2019---formaciones-plana-y-truxe9bol-con-suelo-o-ftb}

\textbf{Referencia.} Quan, L., Fu, C., Si, W., \& Yang, J. (2019).
Numerical study of heat transfer in underground power cable system.
\emph{Energy Procedia, 158}, 5317--5322.
https://doi.org/10.1016/j.egypro.2019.01.636

Quan et al.~resuelven en COMSOL un modelo estacionario 2D de tres cables
enterrados y comparan la geometría, la profundidad y el relleno térmico.
El dominio mide 20 m × 10 m, la superficie se mantiene a 20 °C y las
demás fronteras son adiabáticas.

{\def\LTcaptype{none} % do not increment counter
\begin{longtable}[]{@{}
  >{\raggedright\arraybackslash}p{(\linewidth - 2\tabcolsep) * \real{0.3200}}
  >{\raggedright\arraybackslash}p{(\linewidth - 2\tabcolsep) * \real{0.6800}}@{}}
\toprule\noalign{}
\begin{minipage}[b]{\linewidth}\raggedright
Variable
\end{minipage} & \begin{minipage}[b]{\linewidth}\raggedright
Valor o escenario
\end{minipage} \\
\midrule\noalign{}
\endhead
\bottomrule\noalign{}
\endlastfoot
Configuración & Tres cables en formación plana o trébol. \\
Profundidad & 1.0 a 2.0 m. \\
Separación entre ejes, formación plana & 0.2 a 0.6 m. \\
Radios exteriores conductor/aislamiento/pantalla/cubierta & 24.8, 55.3,
61.7 y 66.8 mm. \\
Corriente del análisis térmico & 1145 A. \\
Criterios térmicos & 65 °C como temperatura óptima y 90 °C como
límite. \\
Medio próximo & Suelo nativo o relleno térmico fluidizado (FTB), ambos
con k dependiente de la temperatura. \\
\end{longtable}
}

{\def\LTcaptype{none} % do not increment counter
\begin{longtable}[]{@{}
  >{\raggedright\arraybackslash}p{(\linewidth - 6\tabcolsep) * \real{0.1324}}
  >{\raggedleft\arraybackslash}p{(\linewidth - 6\tabcolsep) * \real{0.3088}}
  >{\raggedleft\arraybackslash}p{(\linewidth - 6\tabcolsep) * \real{0.2941}}
  >{\raggedleft\arraybackslash}p{(\linewidth - 6\tabcolsep) * \real{0.2647}}@{}}
\toprule\noalign{}
\begin{minipage}[b]{\linewidth}\raggedright
Formación
\end{minipage} & \begin{minipage}[b]{\linewidth}\raggedleft
Ampacidad sin FTB, A
\end{minipage} & \begin{minipage}[b]{\linewidth}\raggedleft
Ampacidad con FTB, A
\end{minipage} & \begin{minipage}[b]{\linewidth}\raggedleft
Incremento con FTB
\end{minipage} \\
\midrule\noalign{}
\endhead
\bottomrule\noalign{}
\endlastfoot
Plana & 1396.4 & 1599.3 & 14.5\% \\
Trébol & 1248.5 & 1500.9 & 20.2\% \\
\end{longtable}
}

\textbf{Aporte a la tesis.} El caso permite verificar la interacción
entre la configuración, la profundidad y la heterogeneidad próxima al
cable. Además, los valores publicados de ampacidad permiten comparar el
efecto del FTB en las formaciones plana y trébol.

\section{Oladunjoye et al.~(2012) - propiedades térmicas medidas in
situ}\label{oladunjoye-et-al.-2012---propiedades-tuxe9rmicas-medidas-in-situ}

\textbf{Referencia.} Oladunjoye, M. A., Sanuade, O. A., \& Olaojo, A. A.
(2012). In situ determination of thermal resistivity of soil: Case study
of Olorunsogo Power Plant, Southwestern Nigeria. \emph{ISRN Civil
Engineering, 2012}, 591450. https://doi.org/10.5402/2012/591450

Oladunjoye et al.~caracterizan el entorno térmico de la central de
turbinas de gas Olorunsogo mediante diez calicatas de aproximadamente
1.5 m y mediciones con un equipo KD2 Pro. Aunque el estudio no simula un
cable específico, proporciona propiedades del suelo para escenarios
heterogéneos.

{\def\LTcaptype{none} % do not increment counter
\begin{longtable}[]{@{}lll@{}}
\toprule\noalign{}
Variable medida & Rango & Promedio reportado \\
\midrule\noalign{}
\endhead
\bottomrule\noalign{}
\endlastfoot
Resistividad térmica & 0.3407 a 0.7188 K m/W & 0.5643 K m/W \\
Conductividad térmica & 1.391 a 2.935 W/(m K) & --- \\
Temperatura del suelo & 28.72 a 35.39 °C & 32.11 °C \\
Humedad óptima & 13.00 a 16.20\% & --- \\
Densidad seca máxima & 1725.05 a 1930.00 kg/m³ & 1855.61 kg/m³ \\
Porosidad & 39.74 a 45.64\% & --- \\
\end{longtable}
}

\textbf{Aporte a la tesis.} Las mediciones de campo permiten
parametrizar \(k(x,y)\) y la temperatura del terreno con variabilidad
espacial observada. Este caso complementa los escenarios construidos con
propiedades bibliográficas o de laboratorio.

\section{Möbius et al.~(2025) - ampacidad con datos climáticos de
Hamburgo}\label{muxf6bius-et-al.-2025---ampacidad-con-datos-climuxe1ticos-de-hamburgo}

\textbf{Referencia.} Möbius, P., Michael, L.-H., Plenz, M., Schräder,
J., \& Gockel, C. (2025). Energy cable ampacity: Impact of seasonal and
climate-related changes. \emph{Renewable and Sustainable Energy Reviews,
212}, 115348. https://doi.org/10.1016/j.rser.2025.115348

El caso aplica IEC 60287 a tres cables unipolares N2XS2Y 1×35/16 de 6/10
kV, instalados en trébol a 0.5 m de profundidad. Usa 32 años de
registros de 30 estaciones del área metropolitana de Hamburgo y un suelo
de referencia con 30\% de arena, 55\% de limo, 15\% de arcilla y
densidad seca de 1.2 t/m³.

{\def\LTcaptype{none} % do not increment counter
\begin{longtable}[]{@{}lrrr@{}}
\toprule\noalign{}
Condición & Temperatura del suelo & Resistividad térmica & Ampacidad \\
\midrule\noalign{}
\endhead
\bottomrule\noalign{}
\endlastfoot
Fría y húmeda & 0 °C & 1.0 K m/W & 245 A \\
Caliente y seca & 20 °C & 2.5 K m/W & 148 A \\
Valor declarado por el fabricante & --- & --- & 187 A \\
\end{longtable}
}

El suelo de referencia tiene punto de marchitez de 12\%, capacidad de
campo de 24\% y saturación máxima de 37\%. Para el análisis de
sensibilidad, los autores variaron la temperatura entre 0 y 45 °C y la
resistividad entre 0.1 y 4.0 K m/W.

\textbf{Aporte a la tesis.} La serie histórica define una instancia
operacional con variación estacional del suelo. La diferencia de 97 A
entre las condiciones extremas permite evaluar el efecto de representar
la temperatura y la resistividad mediante valores variables en lugar de
constantes.

\section{Casos normativos y de verificación de IEC y
CIGRE}\label{casos-normativos-y-de-verificaciuxf3n-de-iec-y-cigre}

Las fuentes normativas cumplen funciones complementarias dentro del
conjunto de verificación. IEC 60287-1-1 aporta las ecuaciones de
ampacidad estacionaria, pérdidas y parámetros térmicos; CIGRE TB 880
desarrolla casos para comprobar su aplicación. En cambio, IEC 60853-3
incorpora un cálculo cíclico completo en los anexos A y B.

{\def\LTcaptype{none} % do not increment counter
\begin{longtable}[]{@{}
  >{\raggedright\arraybackslash}p{(\linewidth - 4\tabcolsep) * \real{0.0968}}
  >{\raggedright\arraybackslash}p{(\linewidth - 4\tabcolsep) * \real{0.4839}}
  >{\raggedright\arraybackslash}p{(\linewidth - 4\tabcolsep) * \real{0.4194}}@{}}
\toprule\noalign{}
\begin{minipage}[b]{\linewidth}\raggedright
Fuente
\end{minipage} & \begin{minipage}[b]{\linewidth}\raggedright
Función en la verificación
\end{minipage} & \begin{minipage}[b]{\linewidth}\raggedright
Uso en el conjunto de casos
\end{minipage} \\
\midrule\noalign{}
\endhead
\bottomrule\noalign{}
\endlastfoot
IEC 60287-1-1 & Aporta formulación y parámetros para el cálculo
estacionario. & Mantener como línea base normativa, no como instancia
autónoma. \\
IEC 60853-3, anexos A y B & Sí; cálculo cíclico con secado parcial y
datos completos. & Incorporar como instancia normativa transitoria. \\
CIGRE TB 880 & Sí; casos para verificar herramientas basadas en IEC. &
Incorporar el caso introductorio 0 y sus variantes. \\
CIGRE TB 963 & Sí; ejemplo analítico y casos FEM. & Incorporar el anillo
analítico y el caso 1 en ductos. \\
\end{longtable}
}

\section{CIGRE TB 880 (2022) - caso introductorio 0 de cable 132
kV}\label{cigre-tb-880-2022---caso-introductorio-0-de-cable-132-kv}

\textbf{Referencia.} CIGRE Working Group B1.56 (2022). \emph{Power cable
rating examples for calculation tool verification}. Technical Brochure
880, caso introductorio 0, pp.~61--125.

El caso base representa tres cables unipolares de 132 kV y 630 mm² Cu,
con aislamiento XLPE, pantalla laminada de aluminio y cubierta PE. Los
cables, de 75.5 mm de diámetro exterior, están en trébol tocándose,
directamente enterrados a 1.0 m hasta el centro del circuito. Se fija
una temperatura del terreno de 20 °C, resistividad térmica de 1.0 K m/W
y límite del conductor de 90 °C.

{\def\LTcaptype{none} % do not increment counter
\begin{longtable}[]{@{}
  >{\raggedright\arraybackslash}p{(\linewidth - 4\tabcolsep) * \real{0.2857}}
  >{\raggedright\arraybackslash}p{(\linewidth - 4\tabcolsep) * \real{0.4545}}
  >{\raggedleft\arraybackslash}p{(\linewidth - 4\tabcolsep) * \real{0.2597}}@{}}
\toprule\noalign{}
\begin{minipage}[b]{\linewidth}\raggedright
Variante del caso base
\end{minipage} & \begin{minipage}[b]{\linewidth}\raggedright
Tratamiento de pérdidas de pantalla
\end{minipage} & \begin{minipage}[b]{\linewidth}\raggedleft
Ampacidad convergida
\end{minipage} \\
\midrule\noalign{}
\endhead
\bottomrule\noalign{}
\endlastfoot
Pantallas unidas en ambos extremos, resultado IEC & Corrientes
circulantes; pérdidas por corrientes parásitas omitidas & 821.78 A \\
Pantallas unidas en un punto & Sin corrientes circulantes; incluye
pérdidas por corrientes parásitas & 886.18 A \\
Pantallas unidas en ambos extremos, recomendación CIGRE & Corrientes
circulantes y parásitas, con factor de reducción & 803.16 A \\
\end{longtable}
}

La familia completa extiende el mismo cable a instalación directa,
ductos HDPE, ductos PVC en concreto, aire con radiación solar y canal
sin relleno. Esta secuencia permite reproducir primero las resistencias,
las pérdidas y la ampacidad del circuito IEC, y luego evaluar los
cambios de geometría, frontera y tratamiento de pérdidas.

\textbf{Aporte a la tesis.} El caso funciona como referencia normativa
de regresión. Aunque no proporciona un campo 2D, permite comprobar la
relación corriente--pérdidas--temperatura y comparar la ampacidad
estimada por PINN o FEM con valores convergidos.

\section{CIGRE TB 963 (2025) - caso 1 FEM 2D con cables en
ductos}\label{cigre-tb-963-2025---caso-1-fem-2d-con-cables-en-ductos}

\textbf{Referencia.} CIGRE Working Group B1.87 (2025). \emph{Finite
element analysis for cable rating calculations}. Technical Brochure 963,
caso 1, pp.~103--121.

La configuración contiene tres cables de 132 kV y 630 mm² Cu en ductos
HDPE llenos de aire, dispuestos horizontalmente dentro de un backfill y
rodeados por suelo y una capa superficial de asfalto. El caso comparte
el diseño de cable de TB 880, pero explicita el dominio y los materiales
necesarios para una referencia FEM 2D.

{\def\LTcaptype{none} % do not increment counter
\begin{longtable}[]{@{}
  >{\raggedright\arraybackslash}p{(\linewidth - 2\tabcolsep) * \real{0.6429}}
  >{\raggedright\arraybackslash}p{(\linewidth - 2\tabcolsep) * \real{0.3571}}@{}}
\toprule\noalign{}
\begin{minipage}[b]{\linewidth}\raggedright
Parámetro
\end{minipage} & \begin{minipage}[b]{\linewidth}\raggedright
Valor
\end{minipage} \\
\midrule\noalign{}
\endhead
\bottomrule\noalign{}
\endlastfoot
Diámetros exteriores conductor / pantallas semiconductoras / XLPE /
pantalla metálica / cubierta & 30.30 / 33.30 / 64.30 / 66.90 / 68.50 /
75.50 mm \\
Profundidad al centro de los ductos & 1.0 m \\
Ducto HDPE & 140 mm exterior; 10.3 mm de espesor \\
Banco de ductos y separación & 1.0 m × 0.5 m; 0.3 m entre ejes \\
Capa de asfalto & 0.1 m \\
Temperaturas & Terreno 20 °C; conductor máximo 90 °C \\
Resistividades térmicas & Suelo 1.0 K m/W; backfill 0.5 K m/W \\
Dominio seleccionado tras sensibilidad & 40 m × 20 m \\
\end{longtable}
}

{\def\LTcaptype{none} % do not increment counter
\begin{longtable}[]{@{}
  >{\raggedright\arraybackslash}p{(\linewidth - 4\tabcolsep) * \real{0.2759}}
  >{\raggedright\arraybackslash}p{(\linewidth - 4\tabcolsep) * \real{0.3793}}
  >{\raggedleft\arraybackslash}p{(\linewidth - 4\tabcolsep) * \real{0.3448}}@{}}
\toprule\noalign{}
\begin{minipage}[b]{\linewidth}\raggedright
Variante
\end{minipage} & \begin{minipage}[b]{\linewidth}\raggedright
Descripción
\end{minipage} & \begin{minipage}[b]{\linewidth}\raggedleft
Ampacidad
\end{minipage} \\
\midrule\noalign{}
\endhead
\bottomrule\noalign{}
\endlastfoot
0 & Modelo base; pantallas con unión multipunto & 615 A \\
1 & Transferencia de calor del aire en ductos modelada explícitamente &
630 A \\
2 & Pantallas con unión en un punto & 973 A \\
3 & Ductos rellenos con bentonita & 700 A \\
4 & Dominios infinitos en los bordes & 615 A \\
5 & Fuentes térmicas definidas según IEC & 617 A \\
6 & Secado del suelo & 587 A \\
7 & Pantalla de alambres de cobre & 654 A \\
8 & Carga variable transitoria; máximo de 85 °C & 615 A \\
\end{longtable}
}

En la variante base, a 615 A las temperaturas de conductor de las tres
fases son 90.0, 83.6 y 82.0 °C; las del aire de los ductos son 71.2,
68.4 y 66.1 °C. El barrido del dominio produce 615.84 A con 200 m × 100
m y 609.85 A con 10 m × 5 m, lo que permite evaluar sensibilidad a la
frontera artificial.

\textbf{Aporte a la tesis.} El caso reúne los componentes físicos del
artefacto propuesto: geometría 2D, varios materiales, fuentes
dependientes de corriente, interfaces y resultados por fase. Por ello,
permite comprobar el campo térmico, el punto caliente, el balance de
energía, la sensibilidad al dominio y la ampacidad.

\section{IEC 60853-3 (2002) - ejemplo cíclico con secado
parcial}\label{iec-60853-3-2002---ejemplo-cuxedclico-con-secado-parcial}

\textbf{Referencia.} International Electrotechnical Commission (2002).
\emph{IEC 60853-3: Calculation of the cyclic and emergency current
rating of cables--Part 3: Cyclic rating factor for cables of all
voltages, with partial drying of the soil}, anexos A y B, pp.~21--29.

El ejemplo considera un cable tripolar de 132 kV, conductor de cobre de
400 mm², aislamiento de papel impregnado en aceite, cubierta de plomo,
refuerzo y cubierta PVC. Como el alcance inicial es estacionario, esta
instancia se reserva para evaluar posteriormente el término temporal y
los efectos del secado.

{\def\LTcaptype{none} % do not increment counter
\begin{longtable}[]{@{}
  >{\raggedright\arraybackslash}p{(\linewidth - 2\tabcolsep) * \real{0.2174}}
  >{\raggedright\arraybackslash}p{(\linewidth - 2\tabcolsep) * \real{0.7826}}@{}}
\toprule\noalign{}
\begin{minipage}[b]{\linewidth}\raggedright
Grupo
\end{minipage} & \begin{minipage}[b]{\linewidth}\raggedright
Datos principales
\end{minipage} \\
\midrule\noalign{}
\endhead
\bottomrule\noalign{}
\endlastfoot
Cable & Diámetro exterior 0.109 m; conductor 22.85 mm; aislamiento 9.5
mm; resistencia AC 61.5 µΩ/m \\
Operación & Ambiente 20 °C; conductor 85 °C; incremento admisible 65 K;
pérdidas dieléctricas 6.03 W/m por cable \\
Circuito térmico & \(T_1=0.835\) K m/W; \(T_3=0.09\) K m/W;
\(T_A=0.278\) K m/W; \(T_B=0.1021\) K m/W \\
Suelo & Profundidad 1.0 m; \(\rho_w=1.0\) y \(\rho_d=2.5\) K m/W;
incremento crítico 30 K; difusividad húmeda \(5\times10^{-7}\) m²/s \\
Ciclo & Factor de pérdidas de 24 h: 0.504; ordenadas de las seis horas
previas: 0.992, 0.728, 0.640, 0.596, 0.593 y 0.796 \\
Resultado & Factor cíclico \(M=1.218\); corregido por secado
\(M_1=1.27\); ampacidad estacionaria con secado 520 A; pico 640 A \\
\end{longtable}
}

\textbf{Aporte a la tesis.} La evaluación inicial mantiene este caso
separado de las referencias estacionarias. Su clasificación como caso
normativo transitorio conserva los parámetros necesarios para una prueba
posterior de la extensión temporal del artefacto.

\section{Pan et al.~(2025) - placa 2D para reconstrucción térmica con
PINN}\label{pan-et-al.-2025---placa-2d-para-reconstrucciuxf3n-tuxe9rmica-con-pinn}

\textbf{Referencia.} Pan, Y., Zhang, K., Ma, N., \& Zhang, J. (2025).
Research on the reconstruction of the temperature field in
two-dimensional steady-state thermal conductivity based on
physics-informed neural networks. \emph{Eng, 6}, 99.
https://doi.org/10.3390/eng6050099

Pan et al.~estudian una placa rectangular \(x\in[0,1]\) m,
\(y\in[0,0.5]\) m, sin fuente interna y con conductividad constante, por
lo que satisface \(\nabla^2T=0\). La temperatura es de 25 °C en el borde
inferior y de 0 °C en los otros tres lados. En el problema inverso, la
condición inferior se reconstruye a partir de datos interiores.

{\def\LTcaptype{none} % do not increment counter
\begin{longtable}[]{@{}
  >{\raggedright\arraybackslash}p{(\linewidth - 2\tabcolsep) * \real{0.6154}}
  >{\raggedright\arraybackslash}p{(\linewidth - 2\tabcolsep) * \real{0.3846}}@{}}
\toprule\noalign{}
\begin{minipage}[b]{\linewidth}\raggedright
Elemento
\end{minipage} & \begin{minipage}[b]{\linewidth}\raggedright
Valor
\end{minipage} \\
\midrule\noalign{}
\endhead
\bottomrule\noalign{}
\endlastfoot
Referencia numérica & ANSYS/Fluent; 10 153 elementos y 10 368 nodos \\
Datos etiquetados & 20 puntos interiores, aproximadamente 0.2\% de los
nodos \\
Entrenamiento reportado & 100 000 iteraciones; diez repeticiones y
selección del mejor resultado \\
Resultado base & Error relativo promedio aproximado de 7.8\% \\
Generalización & Para bordes entre 10 y 40 °C: error relativo promedio
menor de 10\% y error absoluto menor de 1 °C \\
Ponderación adaptativa & Error absoluto máximo reducido a 0.6 °C y error
relativo promedio reducido en unos 2 puntos porcentuales \\
\end{longtable}
}

\textbf{Aporte a la tesis.} Esta instancia 2D permite comprobar la red,
el muestreo y las métricas en una geometría simple que no depende de un
cable. Por tanto, precede a los casos con fuentes cilíndricas,
materiales discontinuos y cálculo de ampacidad.

\section{CIGRE TB 963 (2025) - caso analítico
anular}\label{cigre-tb-963-2025---caso-analuxedtico-anular}

En las pp.~52--55, CIGRE TB 963 define un anillo con radio interior 0.06
m y exterior 20 m. El material tiene resistividad térmica \(\rho=0.8\) K
m/W; el borde exterior está a 20 °C y el interior recibe un flujo de
106.1 W/m², equivalente a 40 W/m. La solución exacta es

\[
T(r)=20-\frac{40(0.8)}{2\pi}\ln\left(\frac{r}{20}\right),
\qquad r=\sqrt{x^2+y^2}.
\]

La temperatura exacta en el radio interior es aproximadamente 49.59 °C.
Aunque la solución es radial, el anillo puede representarse como un
dominio 2D en coordenadas cartesianas. Así, la solución proporciona una
referencia exacta en todo el dominio y no solo en el punto de
temperatura máxima.

\textbf{Aporte a la tesis.} La investigación adopta este anillo como
caso inicial para verificar las derivadas automáticas, las condiciones
de flujo y temperatura, el error de campo y la convergencia. Una vez
superadas estas comprobaciones, la evaluación avanza hacia la geometría
multicapa del cable.

\section{Hahn y Özişik (2012) - rectángulo 2D con frontera
convectiva}\label{hahn-y-uxf6ziux15fik-2012---rectuxe1ngulo-2d-con-frontera-convectiva}

\textbf{Referencia.} Hahn, D. W., \& Özişik, M. N. (2012). \emph{Heat
Conduction} (3.ª ed.), ejemplo 3-3, pp.~92--95. Wiley.
https://doi.org/10.1002/9781118411285

El ejemplo rectangular de Hahn y Özişik incorpora una condición mixta
que complementa las fronteras prescritas de la placa de Pan. En
\(0<x<L\), \(0<y<W\) se resuelve \(\nabla^2T=0\), con \(T=0\) en
\(x=0\), \(x=L\) y \(y=0\), y

\[
-k\frac{\partial T}{\partial y}=h[T(x,W)-T_\infty]
\]

en el borde superior. La solución exacta es

\[
T(x,y)=\sum_{n=1}^{\infty}C_n\sin(\lambda_n x)\sinh(\lambda_n y),
\qquad \lambda_n=\frac{n\pi}{L},
\]

\[
C_n=\frac{hT_\infty\int_0^L\sin(\lambda_nx)\,dx}
{[k\lambda_n\cosh(\lambda_nW)+h\sinh(\lambda_nW)]\int_0^L\sin^2(\lambda_nx)\,dx}.
\]

El caso inicial usa \(L=1\) m, \(W=0.5\) m, \(k=1\) W/(m K), \(h=10\)
W/(m² K) y \(T_\infty=40\) °C. El número de términos de la serie se
registra para mantener el error de truncamiento por debajo del criterio
exigido al artefacto PINN.

\textbf{Aporte a la tesis.} El caso verifica una condición de Robin que
representa la convección equivalente en la superficie del terreno o en
un cable expuesto. Además, permite separar el error del artefacto PINN
del error de truncamiento de la serie.

\section{Fuentes generales de transferencia de
calor}\label{fuentes-generales-de-transferencia-de-calor}

Las fuentes generales de transferencia de calor amplían la verificación
más allá de las instalaciones de cables. Carslaw y Jaeger, Hahn y
Özişik, y Kakac et al.~desarrollan soluciones de conducción en
coordenadas rectangulares y cilíndricas, mientras que Patankar aporta
una referencia para la discretización conservativa y el balance de
flujo.

Dentro del alcance 2D estacionario, el ejemplo rectangular de Hahn y
Özişik se incorpora como instancia explícita porque combina una frontera
convectiva con una solución analítica en serie. Los demás textos
conforman un catálogo para futuras pruebas transitorias, cilíndricas o
tridimensionales.

Xing et al.~(2023) presentan tres problemas anisotrópicos con soluciones
analíticas en dominios 3D. Estos problemas no se incorporan directamente
porque la unidad de análisis es 2D; sin embargo, su estrategia de
definir una solución exacta y deducir la ecuación y sus contornos
orienta la suite manufacturada siguiente.

\section{Suite propuesta de soluciones manufacturadas
2D}\label{suite-propuesta-de-soluciones-manufacturadas-2d}

Las dos instancias manufacturadas se definen a partir de la ecuación de
conducción y del criterio de comparar PINN y FEM con soluciones exactas.
Corresponden a una elaboración propia y se mantienen separadas de los
resultados publicados. En cada caso, la fuente \(Q\) se obtiene al
sustituir una temperatura prescrita en

\[
-\nabla\cdot\left(k(x,y)\nabla T\right)=Q(x,y).
\]

\subsection{MMS-2D-01: conductividad espacialmente
variable}\label{mms-2d-01-conductividad-espacialmente-variable}

En el cuadrado unitario se define

\[
T=T_0+A\sin(\pi x)\sin(\pi y),\qquad
k=k_0(1+\beta x),
\]

con \(T=T_0\) en todo el borde. La fuente compatible es

\[
Q=2k_0(1+\beta x)A\pi^2\sin(\pi x)\sin(\pi y)
-k_0\beta A\pi\cos(\pi x)\sin(\pi y).
\]

El caso comprueba simultáneamente la implementación de \(k(x,y)\), su
gradiente, la fuente y las fronteras de Dirichlet. Los valores iniciales
son \(T_0=20\) °C, \(A=30\) K, \(k_0=1\) W/(m K) y \(\beta=1\).

\subsection{MMS-2D-02: interfaz de dos
materiales}\label{mms-2d-02-interfaz-de-dos-materiales}

El segundo caso define una interfaz vertical \(x=a\), con \(k_1\) a la
izquierda y \(k_2\) a la derecha. La temperatura se expresa como

\[
T=T_0+\phi(x)\sin(\pi y),
\]

\[
\phi(x)=
\begin{cases}
A x/k_1, & x\le a,\\
A\left[a/k_1+(x-a)/k_2\right], & x>a.
\end{cases}
\]

En cada material, \(Q_i=k_i\pi^2\phi(x)\sin(\pi y)\). La construcción
conserva la continuidad de \(T\) y de \(k\,\partial T/\partial x\) en la
interfaz. Los valores \(a=0.5\), \(k_1=0.5\), \(k_2=2.0\) W/(m K),
\(A=10\) W/m y \(T_0=20\) °C definen el contraste inicial de
conductividad.

{\def\LTcaptype{none} % do not increment counter
\begin{longtable}[]{@{}
  >{\raggedright\arraybackslash}p{(\linewidth - 2\tabcolsep) * \real{0.4727}}
  >{\raggedright\arraybackslash}p{(\linewidth - 2\tabcolsep) * \real{0.5273}}@{}}
\toprule\noalign{}
\begin{minipage}[b]{\linewidth}\raggedright
Prueba de aceptación común
\end{minipage} & \begin{minipage}[b]{\linewidth}\raggedright
Criterio que debe registrarse
\end{minipage} \\
\midrule\noalign{}
\endhead
\bottomrule\noalign{}
\endlastfoot
Exactitud de campo & MAE, RMSE, NRMSE y norma relativa \(L_2\) frente a
\(T\) exacta \\
Punto caliente & Error absoluto de \(T_{max}\) y distancia de su
localización \\
Consistencia física & Norma del residuo de PDE y balance entre fuente y
flujo de borde \\
Interfaz & Saltos de temperatura y flujo normal para MMS-2D-02 \\
Robustez & Distribución de errores y convergencia en repeticiones con
semillas declaradas \\
\end{longtable}
}

\textbf{Aporte a la tesis.} La suite evalúa dos condiciones que el
anillo homogéneo no representa: una conductividad variable suave y una
interfaz discontinua. Sus datos se registran como elaboración propia y
se mantienen separados de los casos publicados.

\section{Variables consolidadas de los
casos}\label{variables-consolidadas-de-los-casos}

Los datos extraídos se organizan como una biblioteca de casos
reproducibles del objeto de estudio y de verificación del artefacto. La
tabla siguiente relaciona cada grupo de variables con su fuente y con el
uso previsto en los escenarios de simulación.

{\def\LTcaptype{none} % do not increment counter
\begin{longtable}[]{@{}
  >{\raggedright\arraybackslash}p{(\linewidth - 6\tabcolsep) * \real{0.0954}}
  >{\raggedright\arraybackslash}p{(\linewidth - 6\tabcolsep) * \real{0.4481}}
  >{\raggedright\arraybackslash}p{(\linewidth - 6\tabcolsep) * \real{0.1535}}
  >{\raggedright\arraybackslash}p{(\linewidth - 6\tabcolsep) * \real{0.3029}}@{}}
\toprule\noalign{}
\begin{minipage}[b]{\linewidth}\raggedright
Grupo de datos
\end{minipage} & \begin{minipage}[b]{\linewidth}\raggedright
Variables extraídas
\end{minipage} & \begin{minipage}[b]{\linewidth}\raggedright
Fuentes
\end{minipage} & \begin{minipage}[b]{\linewidth}\raggedright
Uso recomendado
\end{minipage} \\
\midrule\noalign{}
\endhead
\bottomrule\noalign{}
\endlastfoot
Cable multicapa & Diámetros, espesores, materiales, k de
conductor/aislamiento/pantallas/cubierta. & Kim, Al-Dulaimi, Atoccsa,
Ocłoń, Aras, Quan & Definir geometría interna o equivalente térmico del
cable. \\
Instalación & Profundidad H/L, separación S/l/s, formación plana o
trébol, ductos, arreglo de 1/3/6 cables. & Kim, Al-Dulaimi, Atoccsa,
Ocłoń, Aras, Quan, Möbius & Generar dominios 2D reproducibles y variar
proximidad de fuentes. \\
Suelo nativo & k, resistividad térmica, humedad, clasificación,
temperatura y densidad del suelo. & Kim, Khumalo, Aras, Ocłoń,
Oladunjoye, Möbius & Modelar condición homogénea, estratificada, seca o
medida en campo. \\
Bedding/backfill & PAC, NC, arena, SCMB, FTB, SBM, HDPE, conductividad y
dimensiones. & Kim, Al-Dulaimi, Atoccsa, Ocłoń, Quan & Construir
heterogeneidades de alta/baja k y escenarios de mejora térmica. \\
Condiciones de frontera & Temperatura superficial/ambiente, convección,
viento, temperatura inferior constante, fronteras adiabáticas. & Kim,
Al-Dulaimi, Ocłoń, Aras, Quan, Möbius & Definir problemas directos
comparables con FEM/PINN. \\
Resultados térmicos & Tmax, diferencias por posición, campo de
temperatura, criterio 65/90 °C. & Kim, Al-Dulaimi, Ocłoń, Aras, Quan &
Validar salida de campo y localización del punto caliente. \\
Resultados operativos & Ampacidad Imax, reducción por secado, incremento
con backfill y variación estacional. & Khumalo, Atoccsa, Aras, Kim,
Quan, Möbius & Evaluar estimación de Imax y sensibilidad a k y
temperatura del suelo. \\
Verificación normativa & Resistencias, pérdidas, temperatura, factor
cíclico y ampacidad convergida. & CIGRE TB 880; IEC 60853-3 & Probar la
cadena corriente--pérdidas--temperatura y una futura extensión
transitoria. \\
Referencia FEM 2D & Dominio, malla, materiales, fronteras, temperaturas
por fase y sensibilidad. & CIGRE TB 963, caso 1; Pan et al. & Comparar
campo, punto caliente, dominio artificial y reconstrucción inversa. \\
Referencia analítica & Solución exacta de campo en anillo o rectángulo y
condición de Robin. & CIGRE TB 963, ejemplo 4.3; Hahn y Özişik & Medir
error de campo sin incertidumbre de una malla FEM. \\
Verificación manufacturada & \(T(x,y)\) exacta, \(k(x,y)\), fuente
compatible, fronteras e interfaces. & Suite MMS-2D propuesta & Probar
conductividad variable, derivadas automáticas y continuidad
multimaterial. \\
\end{longtable}
}

{\def\LTcaptype{none} % do not increment counter
\begin{longtable}[]{@{}
  >{\raggedright\arraybackslash}p{(\linewidth - 4\tabcolsep) * \real{0.2560}}
  >{\raggedright\arraybackslash}p{(\linewidth - 4\tabcolsep) * \real{0.1280}}
  >{\raggedright\arraybackslash}p{(\linewidth - 4\tabcolsep) * \real{0.6160}}@{}}
\toprule\noalign{}
\begin{minipage}[b]{\linewidth}\raggedright
Nivel de uso
\end{minipage} & \begin{minipage}[b]{\linewidth}\raggedright
Caso recomendado
\end{minipage} & \begin{minipage}[b]{\linewidth}\raggedright
Razón
\end{minipage} \\
\midrule\noalign{}
\endhead
\bottomrule\noalign{}
\endlastfoot
Verificación simple & Aras 2005 & Suelo homogéneo, dominio FEM claro y
comparación IEC/FEM. \\
Heterogeneidad por humedad & Khumalo 2025 & Datos numéricos directos de
resistividad-humedad-ampacidad. \\
Backfill de alta conductividad & Kim 2025 & Compara arena natural, NC,
PAC y CLSM con temperaturas máximas. \\
Escenarios paramétricos & Al-Dulaimi 2024 & Incluye profundidad,
separación, configuración, suelos multicapa y backfills. \\
Optimización de ampacidad & Atoccsa 2024 & Relaciona dimensiones de
backfill con Imax y costo. \\
Backfill localizado y proximidad & Ocłoń 2015 & Dominio 2D con
FTB/SBM/HDPE y fuente central crítica. \\
Geometría y FTB con k(T) & Quan 2019 & Compara formaciones plana/trébol
y cuantifica el aumento de ampacidad. \\
Propiedades in situ & Oladunjoye 2012 & Diez puntos de campo con k,
resistividad, temperatura y propiedades físicas. \\
Variación estacional & Möbius 2025 & Relaciona clima, estado del suelo y
ampacidad con 32 años de observaciones. \\
Regresión normativa IEC & CIGRE TB 880, caso 0 & Presenta parámetros y
ampacidades convergidas para variantes controladas. \\
FEM 2D multimaterial publicado & CIGRE TB 963, caso 1 & Reúne la
geometría multimaterial 2D y las salidas previstas del artefacto. \\
Extensión transitoria & IEC 60853-3, anexos A/B & Conserva cable, ciclo,
suelo húmedo/seco y corriente pico reproducible. \\
PINN 2D antes del cable & Pan 2025 & Aísla aprendizaje, muestreo e
inversión de fronteras en una placa simple. \\
Solución exacta con flujo & CIGRE TB 963, anillo & Entrega \(T(x,y)\)
exacta para todo el dominio. \\
Solución exacta con convección & Hahn y Özişik, ejemplo 3-3 & Verifica
una condición de Robin mediante una serie analítica. \\
Pruebas unitarias multimateriales & MMS-2D-01/02 & Aíslan \(k(x,y)\)
variable y continuidad de temperatura/flujo en interfaces. \\
\end{longtable}
}

\section{Fuentes usadas sin tabla de caso
directo}\label{fuentes-usadas-sin-tabla-de-caso-directo}

Las demás fuentes brindan soporte metodológico o teórico. Una fuente no
se incorpora como instancia cuando carece de un dominio, una ecuación,
condiciones auxiliares o una referencia cuantitativa reproducible dentro
del alcance 2D.

{\def\LTcaptype{none} % do not increment counter
\begin{longtable}[]{@{}
  >{\raggedright\arraybackslash}p{(\linewidth - 4\tabcolsep) * \real{0.1519}}
  >{\raggedright\arraybackslash}p{(\linewidth - 4\tabcolsep) * \real{0.2996}}
  >{\raggedright\arraybackslash}p{(\linewidth - 4\tabcolsep) * \real{0.5485}}@{}}
\toprule\noalign{}
\begin{minipage}[b]{\linewidth}\raggedright
Tipo
\end{minipage} & \begin{minipage}[b]{\linewidth}\raggedright
Ejemplos
\end{minipage} & \begin{minipage}[b]{\linewidth}\raggedright
Motivo
\end{minipage} \\
\midrule\noalign{}
\endhead
\bottomrule\noalign{}
\endlastfoot
Base analítica / normativa histórica & Neher y McGrath (1957) & Aporta
formulación de cálculo de temperatura y capacidad, no un caso moderno de
instalación con datos completos de suelo/backfill. \\
Revisiones de cable rating & Enescu et al.~(2020, 2021); systematic
mapping de DTR; overview térmico & Sirven para justificar variables y
brechas, pero consolidan literatura en vez de reportar un único dataset
de instalación. \\
PINN y SciML & Raissi et al.; Xing et al.; revisiones de conducción
directa/inversa & Aportan método o casos 3D; no todos corresponden a la
unidad de análisis 2D. \\
Transferencia de calor general & Carslaw y Jaeger; Kakac et al.;
Patankar & Proporcionan familias de soluciones y métodos; se
incorporaron solo los casos con ficha reproducible inmediata. \\
Ciencia del diseño & Peffers; Hevner; Gregor & Soporte metodológico de
DSR, sin datos físicos de cable. \\
\end{longtable}
}

\section{Notas de trazabilidad}\label{notas-de-trazabilidad}

\begin{itemize}
\item
  Los valores se extraen de tablas, figuras y texto técnico de los PDF
  almacenados en la biblioteca local de Zotero.
\item
  Cuando un artículo presenta varios resultados equivalentes, se
  priorizan los que describen el cable, la instalación, el suelo o
  relleno, la frontera térmica y la respuesta \(T_{\max}\)/\(I_{\max}\).
\item
  Las tablas organizan los datos necesarios para preparar los casos de
  ejemplo. En una publicación o capítulo posterior, cada valor numérico
  se vincula con su página, tabla o figura de origen.
\item
  Las soluciones manufacturadas se identifican como elaboración propia y
  se separan de los resultados publicados. Sus fuentes y fronteras se
  generan en forma simbólica y se verifican antes del entrenamiento.
\end{itemize}

\end{document}
